%%%%%%%%%%%%%%%%%%%%%%% CHAPTER - 2 %%%%%%%%%%%%%%%%%%%%\\
\chapter{Literature Review}
\label{C2} %%%%%%%%%%%%%%%%%%%%%%%%%%%%
\graphicspath{{images/}{Figures/}{Figures/Chapter-2figs/}}
\noindent\rule{\linewidth}{2pt}
\textit{This Chapter presents a comprehensive critical review of published scholarly literature, foundational methodologies, and contemporary architectural approaches in university placement management systems, rule-based recommendation engines, resume information extraction pipelines, multi-agent web scraping frameworks, and dynamic constraint satisfaction paradigms.}
%%%%%%%%%%%%%%%%%%%%%%%%%%%%%%%%%%%%%%%%%%%%%%%%%%%%%%%%%%%%%%%%%%%%%%%%%%%%%%%%%%%%%%%%%%%%%%%%%%%%%%%%%%%%%%%%%%%%%%%%%%%%%%%%%%%%
\section{Overview of Campus Placement and Internship Systems} \label{S2.1}
The digitization of campus training and placement operations has received extensive attention in computer science education and educational technology literature over the past two decades \cite{ref1,ref2}. Early institutional software systems focused primarily on centralizing student demographic repositories and automating company visit schedules. These legacy platforms typically functioned as relational database wrappers, presenting static job circulars without contextual relevance ranking or automated constraint validation.

With the proliferation of distributed software applications and web-based portals, placement platforms evolved to support multi-stakeholder interactions between students, departmental placement coordinators, and corporate recruiters \cite{ref3}. However, as documented by recent institutional surveys, standard university portals continue to suffer from significant operational limitations. In particular, they treat job listings as monolithic text documents, shifting the burden of verifying eligibility criteria (such as minimum semester GPA, allowed branches, and commencement dates) entirely onto the student applicant.

\begin{figure}[h!]
	\centering
	\includegraphics[scale=0.62]{fig2_1}
	\caption{Structural Representation of Document Vector Processing and Pattern Mapping in Digital Information Retrieval}
	\label{f2.1}
\end{figure}

The diagram illustrated in Figure~\ref{f2.1} portrays the core information retrieval mechanism foundational to document matching, wherein unstructured profile vectors $u(i)$ are mapped against reference feature criteria with residual noise $\eta(i)$ characterized mathematically as:
\begin{equation}
	v(i) = u(i) + \eta(i) \label{e2.1}
\end{equation}
In the context of automated internship matching, reducing the error residual $\eta(i)$ requires enforcing deterministic boundary conditions rather than relying purely on probabilistic approximations.

%%%%%%%%%%%%%%%%%%%%%%%%%%%%%%%%%%%%%%%%%%%%%%%%%%%%%%%%%%%%%%%%%%%%%%%%%%%%%%%%%%%%%%%%%%%%%%%%%%%%%%%%%%%%%%%%%%%%%%%%%%%%%%%%%
\section{Rule-Based and Deterministic Recommendation Algorithms} \label{S2.2}
Recommender systems in industry predominantly employ collaborative filtering (CF), matrix factorization, or deep neural network embeddings \cite{ref4}. While collaborative filtering excels in e-commerce and media streaming domains by leveraging massive historical interaction matrices, it faces acute theoretical and practical hurdles in university internship matching:
\begin{enumerate}[label=\textbf{\roman*)}]
	\item \textbf{The Cold-Start Dilemma:} University cohorts undergo complete turnover every four years. Incoming penultimate-year students possess zero prior internship application logs, rendering collaborative filtering ineffective.
	\item \textbf{Inadmissibility of Stochastic Hallucinations:} In academic placement, eligibility constraints are non-negotiable. An application that falls 0.05 points below a company's strict CGPA cutoff is rejected instantly by corporate automated screening systems. Probabilistic machine learning models that generate recommendations based on latent similarity can easily recommend ineligible roles.
	\item \textbf{Requirement for Decision Explainability:} Students require deterministic guarantees explaining why a role was placed at Rank 1 versus Rank 2, and exact citations of violated constraints for rejected roles.
\end{enumerate}

To circumvent these vulnerabilities, researchers have advocated for hybrid and rule-based knowledge-driven recommender systems \cite{ref5}. Deterministic engines operate on explicit condition-action rules and discrete mathematical scoring functions. By decoupling hard constraint filtering from soft relevance weighting, deterministic systems provide 100\% reproducible, auditable, and transparent outcomes.

%%%%%%%%%%%%%%%%%%%%%%%%%%%%%%%%%%%%%%%%%%%%%%%%%%%%%%%%%%%%%%%%%%%%%%%%%%%%%%%%%%%%%%%%%%%%%%%%%%%%%%%%%%%%%%%%%%%%%%%%%%%%%%%%%
\section{Automated Resume Parsing and Information Extraction Techniques} \label{S2.3}
The ingestion of candidate profiles from unstructured digital resumes represents a classic challenge in Natural Language Processing (NLP) and Named Entity Recognition (NER) \cite{ref1,ref3}. Modern resumes are distributed predominantly in Portable Document Format (PDF) with multi-column layouts, diverse typographical hierarchies, and idiosyncratic stylistic conventions.

Early research explored optical character recognition (OCR) and Hidden Markov Models (HMM) for document segment classification. Contemporary production frameworks, however, divide the extraction pipeline into two distinct phases:
\begin{itemize}
	\item \textbf{Stream Text Extraction:} Tools such as \texttt{pypdf} extract continuous stream text directly from the underlying PDF document object model, preserving sequential word tokens without rasterizing pages into computationally expensive pixel bitmaps.
	\item \textbf{Lexical Dictionary and Regex Pattern Matching:} Extracted text is normalized and queried against comprehensive technical skill dictionaries and regular expression grammars. Specific lexical patterns identify degree titles (e.g., \texttt{B.Tech CSE}, \texttt{B.E}, \texttt{MCA}), cohort years (e.g., \texttt{3rd Year}, \texttt{Year III}, \texttt{6th Semester}), and numeric CGPA scales (including normalized conversion from 4.0 GPA to 10.0 CGPA formats).
\end{itemize}
This deterministic lexical approach delivers sub-second parsing latency on client devices while eliminating external cloud API dependencies.

%%%%%%%%%%%%%%%%%%%%%%%%%%%%%%%%%%%%%%%%%%%%%%%%%%%%%%%%%%%%%%%%%%%%%%%%%%%%%%%%%%%%%%%%%%%%%%%%%%%%%%%%%%%%%%%%%%%%%%%%%%%%%%%%%
\section{Multi-Agent Web Scraping and External Opportunity Aggregation} \label{S2.4}
While on-campus placement cells provide dedicated internal drives, a substantial proportion of elite internship opportunities are published directly on corporate career portals. Manual monitoring of dozens of external portals is prohibitive for individual students. 

Recent literature has highlighted the efficacy of multi-agent web scraping frameworks \cite{ref2,ref4}. In an autonomous multi-agent crawling architecture:
\begin{itemize}
	\item \textbf{Harvester Agents:} Dispatch specialized HTTP requests targeting endpoint query structures across enterprise recruitment architectures (e.g., Workday, Oracle Taleo, Greenhouse, Lever, and proprietary Google/Amazon careers APIs).
	\item \textbf{Sanitization and Parsing Agents:} Strip redundant Document Object Model (DOM) boilerplate, navigation menus, and JavaScript tracking snippets, isolating the central job description text.
	\item \textbf{Semantic Extraction Agents:} Employ fast inference language models (such as LLaMA-3.1-8B-Instant via Groq API) or rule-based fallback NLP to transform raw listing text into standardized structured JSON objects conforming to the system's schema.
\end{itemize}

%%%%%%%%%%%%%%%%%%%%%%%%%%%%%%%%%%%%%%%%%%%%%%%%%%%%%%%%%%%%%%%%%%%%%%%%%%%%%%%%%%%%%%%%%%%%%%%%%%%%%%%%%%%%%%%%%%%%%%%%%%%%%%%%%
\section{Dynamic Constraint Handling and Reassessment in Recommender Systems} \label{S2.5}
A critical yet frequently neglected challenge in recommender systems is \textit{state drift}---the scenario where opportunity availability or candidate eligibility changes dynamically during an active selection cycle \cite{ref5}. In campus hiring:
\begin{itemize}
	\item A company reaches its application quota and marks an internship status as "Closed".
	\item High application volumes cause a recruiter to raise the minimum CGPA cutoff from 8.0 to 8.5.
	\item A candidate's semester examination schedule is finalized, restricting their available commencement date.
\end{itemize}

Standard recommender systems silently drop newly ineligible items from the output feed, leaving the user confused as to why their previously favored opportunity vanished. Scholarly investigations in explainable AI (XAI) emphasize that systems must maintain an explicit \textit{Transition Matrix}. By detailing previously held ranks, updated statuses, agent decisions ("Remove" vs. "Reconsider"), and plain-language justifications, the system establishes complete transparency and trust with student stakeholders.

%%%%%%%%%%%%%%%%%%%%%%%%%%%%%%%%%%%%%%%%%%%%%%%%%%%%%%%%%%%%%%%%%%%%%%%%%%%%%%%%%%%%%%%%%%%%%%%%%%%%%%%%%%%%%%%%%%%%%%%%%%%%%%%%%
\section{Identification of Research Gaps} \label{S2.6}
A comparative analysis of existing commercial platforms (e.g., Internshala, LinkedIn, Glassdoor, Handshake) and academic prototypes reveals four critical research and architectural gaps:
\begin{enumerate}[label=\textbf{\arabic*.}]
	\item \textbf{Absence of Strict Multi-Constraint Academic Filtering:} Generic commercial platforms allow students to apply indiscriminately, resulting in high rejection rates due to hidden institutional prerequisites (e.g., university duration caps).
	\item \textbf{Opaque Scoring and Zero Disqualification Feedback:} Candidates who fail to receive recommendations are rarely provided with actionable explanations detailing missing skills or specific cutoff failures.
	\item \textbf{Lack of Dynamic Shortlist Reassessment Tables:} No existing campus platform provides an autonomous transition table that audits disqualified roles and promotes replacement alternatives when conditions shift.
	\item \textbf{Disjointed Candidate Lifecycles:} Most platforms isolate recommendation from preparation, offering no native bridge to What-If sensitivity testing, Kanban application tracking, or automated technical mock interviews.
\end{enumerate}

%%%%%%%%%%%%%%%%%%%%%%%%%%%%%%%%%%%%%%%%%%%%%%%%%%%%%%%%%%%%%%%%%%%%%%%%%%%%%%%%%%%%%%%%%%%%%%%%%%%%%%%%%%%%%%%%%%%%%%%%%%%%%%%%%
\section{Summary} \label{S2.7}
This Chapter reviewed the historical evolution and current technological state of campus placement systems, recommender algorithms, resume information extraction, and web scraping agents. The critical limitations of probabilistic black-box models in academic environments establish a compelling imperative for the deterministic, explainable, and multi-featured architecture realized in the Campus Agent platform.
